\section{模型检验}

\subsection{约束与数值可行性检验}

首先检验问题一最优方案是否满足各项运行约束。全天储电量保持在
$1200\sim10800$ kWh范围内，单个$10\,\mathrm{min}$时段最大充电量为
$833.333$ kWh，最大放电量为$715.653$ kWh，均未超过题目规定的充放电功率限制。
全天不存在同时充放电，日末储电量恢复至$6000$ kWh，且未发生弃光。

进一步对问题二和问题三的主要策略进行数值校验。
完全信息基准与Dynamic在统计期334天内均满足储电量边界、充放电功率及互斥约束，
储电量能够跨日连续传递。完全信息基准的最大电量平衡残差约为
$8.00\times10^{-9}$ kWh，储电量递推最大残差约为
$9.26\times10^{-8}$ kWh；Dynamic实际结算中的最大电量平衡残差约为
$4.55\times10^{-13}$ kWh，弃光量与外部电网购电溢余量分解的最大日汇总误差约为
$8.19\times10^{-12}$ kWh。

问题三进一步检验储电量递推、实际电量平衡、紧急购电量及费用计算，
结果见表~\ref{tab:p3_check}。

\begin{table}[H]
\centering
\caption{问题三模型求解结果检验}
\label{tab:p3_check}
\begin{tabular}{lc}
\toprule
检验项目 & 最大误差或数量 \\
\midrule
储电量递推残差 & $9.09\times10^{-13}\,\mathrm{kWh}$ \\
实际电量平衡残差 & $4.55\times10^{-13}\,\mathrm{kWh}$ \\
紧急购电量计算残差 & $4.55\times10^{-13}\,\mathrm{kWh}$ \\
每日费用计算误差 & $2.91\times10^{-11}$\,元 \\
同时充放电时段数量 & $0$ \\
储电量越界天数 & $0$ \\
\bottomrule
\end{tabular}
\end{table}

对问题四中波动电价下的问题二主要策略和问题三主要策略进行同样检验，
统计期均为2025年2月1日至12月31日，共334天，结果见表~\ref{tab:p4_numerical_check}。

\begin{table}[H]
\centering
\caption{问题四主要策略的模型求解检验结果}
\label{tab:p4_numerical_check}
\begin{tabular}{lcc}
\toprule
检验指标 & 问题二主要策略 & 问题三主要策略 \\
\midrule
统计天数 & 334 & 334 \\
违反储电量上下限的天数 & 0 & 0 \\
最大电量平衡残差（kWh） & $4.55\times10^{-13}$ & $4.55\times10^{-13}$ \\
储电量最小值（kWh） & 1200 & 1200 \\
储电量最大值（kWh） & 10800 & 10800 \\
最大单时段充电量（kWh） & 833.333 & 833.333 \\
最大单时段放电量（kWh） & 833.333 & 833.333 \\
同时充放电时段数 & 0 & 0 \\
\bottomrule
\end{tabular}
\end{table}

两种主要策略在统计期内均未出现储电量越界或同时充放电，
单时段充放电量均满足功率限制，最大电量平衡残差为
$4.55\times10^{-13}$ kWh，属于数值计算误差范围。


\subsection{求解结果交叉检验}

问题一采用PuLP调用CBC求解混合整数线性规划模型。
为检验最优解的稳定性，去除充放电互斥的0--1变量后，将模型转化为线性规划问题，
并采用SciPy/HiGHS独立求解。两种求解方式得到的最优目标值一致，
且线性规划解中未出现同一时段同时充电和放电，
说明充放电互斥约束未改变问题一的最优目标值。


\subsection{Dynamic信息非泄露检验}

为检验Dynamic日前决策是否使用未来信息，
对目标日期之后尚未发生的负载和光伏发电量数据进行随机扰动，并重新执行模型。
扰动后，目标日期之前的历史数据选择、日前预测和动态安全冗余决策均保持不变。
因此，第$d$天的购电计划仅依赖该决策时点之前已经获得的信息，
目标日及未来实际数据仅用于计划执行后的结算。